%----------------------------------------------------------------------------------------
% VOORWAARDELIJKE KANS. Alle opgaven zijn eigen werk.
%
% Opbouw: Start, leerstof (theorie, werkwijze of uitleg, voorbeeld, korte oefening), en op een
% nieuwe pagina de Opdrachten (door elkaar, stijgende moeilijkheid), Koppel en Test jezelf.
% Labels van dit hoofdstuk beginnen met h2-.
%----------------------------------------------------------------------------------------

\mychapter{2}{Voorwaardelijke kans}

\lkr{Leerplan B+S': LPD 35, met de wenken over voorwaardelijke en marginale kansen in kruistabellen en boomdiagrammen, de wet van de totale kans en de wet van Bayes. De wet van Bayes wordt eerst met aantallen (een kruistabel van 10\,000 personen) uitgelegd en pas daarna met de formule: onderzoek toont dat dat de klassieke denkfout bij medische tests sterk vermindert.}

Is de kans dat het morgen regent dezelfde als de kans dat het morgen regent \textit{als het vandaag regent}? Meestal niet. Extra informatie verandert een kans. In dit hoofdstuk leer je met zulke voorwaardelijke kansen rekenen, met kruistabellen en boomdiagrammen. Op het einde begrijp je waarom een positieve medische test lang niet altijd betekent dat je ziek bent.

\begin{oefeningen}[Start]
\oef{h2-start-1}

\oef{h2-verken}
\end{oefeningen}

\lkr{Laat oefening \ref{oef:h2-verken} maken vóór de uitleg. Het verschil tussen b en c, de kans dat een fietser sport en de kans dat een sporter fietst, is de kern van de voorwaardelijke kans: de voorwaarde bepaalt door welk totaal je deelt.}

%========================================================================================
\section{Kruistabellen}

\subsection{Voorwaardelijke kans}

Een \textbf{kruistabel} telt een groep volgens twee kenmerken tegelijk, zoals de tabel uit de Verken. De totalen staan in de marge: daarom heten de kansen daaruit \textbf{marginale kansen}. Weet je al dat een persoon een kenmerk heeft, bijvoorbeeld dat hij met de fiets komt, dan kijk je alleen nog naar die rij of kolom.

\naast[3.4cm]{Zo werkt het ook algemeen. Weet je dat $B$ is opgetreden, dan liggen de mogelijke uitkomsten in $B$: de uitkomstenverzameling krimpt van $\Omega$ tot $B$. Van $A$ blijft alleen het deel in $B$ over, dus $A \cap B$. Zijn alle uitkomsten even waarschijnlijk, dan is de nieuwe kans op $A$ dus $\frac{\#(A \cap B)}{\#B}$. Deel teller en noemer door $\#\Omega$, dan staan er gewone kansen:
\[\frac{\#(A \cap B)}{\#B} = \frac{\#(A \cap B) / \#\Omega}{\#B / \#\Omega} = \frac{P(A \cap B)}{P(B)} .\]}{\begin{tikzpicture}[font=\small,scale=.85]
  \draw[grijs] (-1.3,-1.1) rectangle (2.5,1.1); \node[grijs,anchor=north east] at (2.5,1.1) {$\Omega$};
  \fill[accent!12] (1.15,0) circle (.85);
  \begin{scope}\clip (1.15,0) circle (.85); \fill[accent!45] (0,0) circle (.85); \end{scope}
  \draw[grijs,dashed] (0,0) circle (.85); \draw[accent,thick] (1.15,0) circle (.85);
  \node[grijs] at (-.4,0) {$A$}; \node[text=accent] at (1.6,0) {$B$};
\end{tikzpicture}}

Die formule nemen we als definitie, ook als de uitkomsten niet even waarschijnlijk zijn.

\begin{theorie}
De \textbf{\textcolor{accent}{voorwaardelijke kans}} op $A$ als $B$ gegeven is, met $P(B) > 0$, is

\vspace{1mm}
\qquad $\formule{P(A \gegeven B) = \dfrac{P(A \cap B)}{P(B)}}$
\end{theorie}

Je leest $P(A \gegeven B)$ als \textit{de kans op $A$ als $B$}. De voorwaarde $P(B) > 0$ is nodig, want je deelt door $P(B)$. In een kruistabel deel je dus door het totaal van de rij of kolom van $B$.

\voorbeeld{}{Een bibliotheek telde 500 ontleningen, naar soort boek en leeftijd van de lener. Wat is de kans dat een ontlening een strip is? Dat een ontlening door een jongere een strip is? Dat een strip ontleend werd door een jongere?\label{vb:h2-kruistabel}}{%
\uitleg{Bij de eerste vraag deel je door het totaal, bij de tweede door het aantal ontleningen door jongeren, bij de derde door het aantal strips.}
\figuren{\renewcommand{\arraystretch}{1.25}\begin{tabular}{l|cc|c}
 & strip & roman & totaal \\ \hline
jongere & 96 & 104 & 200 \\
volwassene & 54 & 246 & 300 \\ \hline
totaal & 150 & 350 & 500
\end{tabular}}
\begin{stappen}
P(S) = \frac{150}{500} = 0{,}30 & Een marginale kans: het totaal van de kolom strip. \\[3ex]
P(S \gegeven J) = \frac{96}{200} = 0{,}48 & Alleen de rij jongere telt mee. \\[3ex]
P(J \gegeven S) = \frac{96}{150} = 0{,}64 & Alleen de kolom strip telt mee.
\end{stappen}}

Let op: $P(S \gegeven J)$ en $P(J \gegeven S)$ zijn verschillende kansen. De teller is dezelfde, de noemer niet. Die twee door elkaar halen is een van de bekendste denkfouten in de kansrekening: de kans dat een strip door een jongere ontleend is (64~\%), is niet de kans dat een jongere een strip ontleent (48~\%).

\begin{oefeningen}
\needspace{12\baselineskip}
\oef{h2-kruistabel}
\end{oefeningen}

\needspace{16\baselineskip}
\subsection{Rekenen met de definitie}

Zonder tabel reken je rechtstreeks met de formule. Vaak moet je daarvoor eerst $P(A \cap B)$ vinden met een rekenregel uit hoofdstuk 1.

\voorbeeld{}{Gegeven: $P(A) = 0{,}6$, $P(B) = 0{,}5$ en $P(A \cup B) = 0{,}8$. Bereken $P(A \gegeven B)$ en $P(B \gegeven A)$.\label{vb:h2-formule}}{%
\uitleg{De definitie vraagt $P(A \cap B)$. Die volgt uit de algemene somregel, als we ze omkeren.}
\begin{stappenbreed}
0{,}8 = 0{,}6 + 0{,}5 - P(A \cap B) & De algemene somregel. \\[1ex]
P(A \cap B) = 0{,}3 & Opgelost naar de onbekende. \\[1ex]
P(A \gegeven B) = \frac{0{,}3}{0{,}5} = 0{,}6 \qquad P(B \gegeven A) = \frac{0{,}3}{0{,}6} = 0{,}5 & Dezelfde teller, maar een andere noemer: deel door de kans op wat gegeven is.
\end{stappenbreed}}

\begin{oefeningen}
\oef{h2-formule}
\end{oefeningen}

%========================================================================================
\section{Boomdiagrammen}

\subsection{De productregel}

Bij twee trekkingen na elkaar ken je vaak de kans op de eerste, en de kans op de tweede \textit{als je weet wat de eerste was}. Uit die twee kansen wil je de kans halen dat beide gebeuren. De definitie van de voorwaardelijke kans geeft dat meteen.

\begin{theorie}
\begin{stelling}{Productregel}
Voor gebeurtenissen $A$ en $B$ met $P(A) > 0$ is \quad $\formule{P(A \cap B) = P(A) \cdot P(B \gegeven A)}$
\end{stelling}

\begin{bewijs}
Volgens de definitie, met de rollen van $A$ en $B$ omgewisseld, is
\bewijsregel{P(B \gegeven A) = \frac{P(B \cap A)}{P(A)} = \frac{P(A \cap B)}{P(A)} .}
Vermenigvuldig beide leden met $P(A)$.
\end{bewijs}
\end{theorie}

In woorden: de kans dat eerst $A$ en dan $B$ gebeurt, is de kans op $A$ maal de kans op $B$ \textit{als $A$ al gebeurd is}. Bij drie stappen ga je zo verder: $P(A \cap B \cap C) = P(A) \cdot P(B \gegeven A) \cdot P(C \gegeven A \cap B)$.

\voorbeeld{}{Je trekt twee kaarten na elkaar uit een spel van 52, zonder de eerste terug te leggen. Wat is de kans op twee azen?\label{vb:h2-product}}{%
\uitleg{Na de eerste trekking verandert het spel: als de eerste kaart een aas is, blijven er 3 azen over in 51 kaarten.}
\begin{stappenbreed}
P(A_1) = \frac{4}{52} & $A_1$: de eerste kaart is een aas. \\[2ex]
P(A_2 \gegeven A_1) = \frac{3}{51} & $A_2$: de tweede is een aas, als de eerste dat al was. \\[2ex]
P(A_1 \cap A_2) = \frac{4}{52} \cdot \frac{3}{51} = \frac{12}{2652} = \frac{1}{221} & De productregel: ongeveer $0{,}45$~\%.
\end{stappenbreed}}

\begin{oefeningen}
\oef{h2-product}
\end{oefeningen}

\needspace{16\baselineskip}
\subsection{Rekenen met een boomdiagram}

Een \textbf{boomdiagram} zet de stappen van een experiment na elkaar. Op elke tak staat de voorwaardelijke kans, gegeven alles wat ervoor kwam. Twee regels volstaan dan:
\begin{itemize}
\item de kans op een pad is het product van de kansen langs het pad: dat is de productregel;
\item leiden meerdere paden naar wat je zoekt, dan tel je hun kansen op: verschillende paden zijn onverenigbaar, dus dat is de somregel.
\end{itemize}

\voorbeeld{}{In een vaas zitten 5 rode en 3 witte knikkers. Je trekt er twee na elkaar, zonder de eerste terug te leggen. Wat is de kans op twee rode? En op één van elke kleur?\label{vb:h2-boom}}{%
\uitleg{Na de eerste trekking zitten er nog 7 knikkers in de vaas, en de kansen bij de tweede trekking hangen af van de eerste. Die voorwaardelijke kansen zetten we op de takken.}
\figuren{\begin{tikzpicture}[x=2.2cm,y=.62cm]
  \node[blad] (s) at (0,0) {start};
  \node[blad] (r) at (1,1.8) {$R$}; \node[blad] (w) at (1,-1.8) {$W$};
  \draw[tak] (s) -- node[kanslabel,above left] {$\frac58$} (r); \draw[tak] (s) -- node[kanslabel,below left] {$\frac38$} (w);
  \node[blad] (rr) at (2,2.7) {$R$}; \node[blad] (rw) at (2,.9) {$W$};
  \node[blad] (wr) at (2,-.9) {$R$}; \node[blad] (ww) at (2,-2.7) {$W$};
  \draw[tak] (r) -- node[kanslabel,above] {$\frac47$} (rr); \draw[tak] (r) -- node[kanslabel,below] {$\frac37$} (rw);
  \draw[tak] (w) -- node[kanslabel,above] {$\frac57$} (wr); \draw[tak] (w) -- node[kanslabel,below] {$\frac27$} (ww);
  \foreach \n/\t in {rr/{RR: \tfrac{20}{56}},rw/{RW: \tfrac{15}{56}},wr/{WR: \tfrac{15}{56}},ww/{WW: \tfrac{6}{56}}}{\node[anchor=west,font=\small] at ($(\n)+(.2,0)$) {$\t$};}
\end{tikzpicture}}
\begin{stappenbreed}
P(RR) = \frac58 \cdot \frac47 = \frac{20}{56} = \frac{5}{14} & Het bovenste pad. \\[3ex]
P(RW) + P(WR) = \frac58 \cdot \frac37 + \frac38 \cdot \frac57 = \frac{30}{56} = \frac{15}{28} & Twee paden: optellen.
\end{stappenbreed}}

De vier eindkansen zijn samen 1: $\frac{20 + 15 + 15 + 6}{56} = 1$. Dat moet, want de vier paden zijn onverenigbaar en samen vormen ze alle mogelijkheden. Het is een goede controle.

\begin{werkwijze}
Teken een boomdiagram met op elke tak de kans, gegeven wat ervoor gebeurde. Vermenigvuldig langs een pad. Tel de paden op die bij de gevraagde gebeurtenis horen. Controleer dat alle eindkansen samen 1 zijn.
\end{werkwijze}

\begin{oefeningen}
\oef{h2-boom}

\oef{h2-boom-2}
\end{oefeningen}

%========================================================================================
\section{Onafhankelijke gebeurtenissen}

\subsection{Definitie}

Leg je de eerste knikker terug, dan verandert de tweede trekking niet door de eerste. De productregel wordt dan eenvoudiger: $P(B \gegeven A)$ is gewoon $P(B)$. Zulke gebeurtenissen heten onafhankelijk. We definiëren dat met de productformule, omdat die ook werkt als een kans 0 is.

\begin{theorie}
Twee gebeurtenissen $A$ en $B$ zijn \textbf{\textcolor{accent}{onafhankelijk}} als \quad $\formule{P(A \cap B) = P(A) \cdot P(B)}$

\begin{stelling}{Onafhankelijkheid en voorwaardelijke kans}
Is $P(B) > 0$, dan zijn $A$ en $B$ onafhankelijk als en slechts als \quad $\formule{P(A \gegeven B) = P(A)}$
\end{stelling}

\begin{bewijs}
Zijn $A$ en $B$ onafhankelijk, dan is
\bewijsregel{P(A \gegeven B) = \frac{P(A \cap B)}{P(B)} = \frac{P(A) \cdot P(B)}{P(B)} = P(A) .}
Omgekeerd: is $P(A \gegeven B) = P(A)$, dan geeft de definitie $\frac{P(A \cap B)}{P(B)} = P(A)$, en na vermenigvuldigen met $P(B)$ staat er $P(A \cap B) = P(A) \cdot P(B)$.
\end{bewijs}
\end{theorie}

Onafhankelijk betekent dus: weten dat $B$ gebeurd is, verandert niets aan de kans op $A$. De definitie is symmetrisch in $A$ en $B$, dus dan verandert weten dat $A$ gebeurd is ook niets aan de kans op $B$.

\voorbeeld{}{Je gooit met twee dobbelstenen. Is \textit{de eerste dobbelsteen is even} ($A$) onafhankelijk van \textit{de som is 7} ($B$)?\label{vb:h2-onafhankelijk}}{%
\uitleg{We tellen in het rooster en vergelijken de kans op beide met het product van de kansen.}
\begin{stappenbreed}
P(A) = \frac{18}{36} = \frac12, \quad P(B) = \frac{6}{36} = \frac16 & Drie even waarden voor de eerste, maal zes voor de tweede; de diagonaal van som 7. \\[3ex]
P(A \cap B) = \frac{3}{36} = \frac{1}{12} & $(2, 5)$, $(4, 3)$ en $(6, 1)$. \\[3ex]
P(A) \cdot P(B) = \frac12 \cdot \frac16 = \frac{1}{12} & Gelijk: onafhankelijk.
\end{stappenbreed}
Welke eerste worp je ook hebt, er is telkens precies één tweede worp die de som 7 maakt. De kans op som 7 blijft dus $\frac16$.}

Onafhankelijkheid zie je niet altijd aan de situatie: je moet ze nagaan met de kansen.

\voorbeeld{}{In de Verken komt de helft van de 200 leerlingen met de fiets ($F$), en sport 45~\% ($S$). Van de fietsers sport 60~\%. Zijn $S$ en $F$ onafhankelijk?\label{vb:h2-afhankelijk}}{%
\uitleg{We vergelijken de voorwaardelijke kans met de gewone kans, zoals de stelling toelaat.}
\begin{stappenbreed}
P(S \gegeven F) = 0{,}60, \quad P(S) = 0{,}45 & Bij de fietsers sporten meer leerlingen dan gemiddeld. \\[1ex]
P(S \gegeven F) \neq P(S) & Dus afhankelijk. Controle met de definitie: $P(S \cap F) = 0{,}30$, maar $P(S) \cdot P(F) = 0{,}225$.
\end{stappenbreed}}

\begin{oefeningen}
\oef{h2-onafhankelijk}
\end{oefeningen}

\needspace{16\baselineskip}
\subsection{Herhaalde experimenten}

Gooi je een munt of een dobbelsteen meerdere keren, of trek je met teruglegging, dan beïnvloeden de herhalingen elkaar niet: ze zijn onafhankelijk. In een boomdiagram staan dan bij elke stap dezelfde kansen, en je vermenigvuldigt ze gewoon. De kans op drie keer zes is $\bigl(\frac16\bigr)^3$.

\voorbeeld{}{Een boogschutter raakt de roos met kans $0{,}3$ per pijl, onafhankelijk van de vorige pijlen. Ze schiet drie pijlen. Wat is de kans dat ze de roos minstens één keer raakt?\label{vb:h2-herhaling}}{%
\uitleg{\textit{Minstens één keer} heeft veel paden. Het tegendeel, \textit{drie keer mis}, is maar één pad.}
\begin{stappenbreed}
P(\text{drie keer mis}) = 0{,}7 \cdot 0{,}7 \cdot 0{,}7 = 0{,}343 & Onafhankelijke pijlen: vermenigvuldig de kans op mis. \\[1ex]
P(\text{minstens één keer raak}) = 1 - 0{,}343 = 0{,}657 & De complementregel.
\end{stappenbreed}}

\begin{oefeningen}
\oef{h2-onafhankelijk-2}
\end{oefeningen}

\needspace{16\baselineskip}
\subsection{Onafhankelijk of onverenigbaar?}

De twee woorden klinken gelijkaardig, maar ze betekenen iets heel anders. \textit{Onverenigbaar} gaat over de gebeurtenissen zelf: ze kunnen niet samen optreden, $A \cap B = \emptyset$. \textit{Onafhankelijk} gaat over kansen: de ene verandert niets aan de kans op de andere. Onverenigbare gebeurtenissen met een positieve kans zijn juist sterk afhankelijk: gebeurt de ene, dan is de kans op de andere 0. Dan is $P(A \cap B) = 0$, terwijl $P(A) \cdot P(B) > 0$.

\voorbeeld{}{Je gooit met een dobbelsteen. Noem $A$ \textit{even}, $B$ \textit{oneven} en $C$ \textit{hoogstens 2}. Zijn $A$ en $B$ onverenigbaar? Onafhankelijk? En $A$ en $C$?\label{vb:h2-onverenigbaar}}{%
\uitleg{Voor onverenigbaar kijken we naar de doorsnede, voor onafhankelijk vergelijken we kansen.}
\begin{stappenbreed}
A \cap B = \emptyset & Onverenigbaar. \\[1ex]
P(A \cap B) = 0 \neq \frac12 \cdot \frac12 & Dus niet onafhankelijk: weet je dat het even is, dan is oneven onmogelijk. \\[1ex]
A \cap C = \{2\}, \quad P(A \cap C) = \frac16 & Niet leeg: niet onverenigbaar. \\[1ex]
P(A) \cdot P(C) = \frac12 \cdot \frac26 = \frac16 & Wel onafhankelijk: van $\{1, 2\}$ is de helft even, net als van $\Omega$.
\end{stappenbreed}}

\begin{oefeningen}
\oef{h2-onverenigbaar}
\end{oefeningen}

%========================================================================================
\section{Totale kans en de wet van Bayes}

\subsection{Een medische test}

Een sneltest voor een virus is positief bij 95~\% van de besmette personen. Bij niet-besmette personen geeft hij in 4~\% van de gevallen toch een positief resultaat: een \textit{vals positief}. In een stad is 2~\% van de inwoners besmet. Iemand test positief. Hoe groot is de kans dat die persoon besmet is? Veel mensen gokken 95~\%. Een tabel met aantallen toont dat het veel minder is.

\voorbeeld{}{Bereken de kans dat iemand die positief test, besmet is. Denk aan 10\,000 inwoners.\label{vb:h2-bayes}}{%
\uitleg{We zetten de percentages om in aantallen voor 10\,000 inwoners, eerst per rij (besmet of niet), dan per kolom (positief of negatief). Dan lezen we af hoeveel van de positieven echt besmet zijn.}
\figuren{\renewcommand{\arraystretch}{1.25}\begin{tabular}{l|cc|c}
 & positief & negatief & totaal \\ \hline
besmet & 190 & 10 & 200 \\
niet besmet & 392 & 9408 & 9800 \\ \hline
totaal & 582 & 9418 & 10\,000
\end{tabular}}
\begin{stappenbreed}
200 \text{ besmet, waarvan } 0{,}95 \cdot 200 = 190 \text{ positief} & 2~\% van 10\,000, en 95~\% daarvan. \\[1ex]
9800 \text{ niet besmet, waarvan } 0{,}04 \cdot 9800 = 392 \text{ positief} & De valse positieven. \\[1ex]
P(\text{besmet} \gegeven \text{positief}) = \frac{190}{582} \approx 0{,}33 & Maar een derde van de positieven is besmet.
\end{stappenbreed}}

Er zijn veel meer valse positieven dan echte, omdat er zoveel meer niet-besmette mensen zijn. Daarom laat een arts een positieve sneltest bevestigen met een betere test. Dezelfde berekening kun je met kansen doen in plaats van aantallen. Dat levert twee formules op: één voor het totaal van een kolom, en één voor de gevraagde voorwaardelijke kans.

\needspace{16\baselineskip}
\subsection{De wet van de totale kans}

In de tabel kwam het aantal positieven uit twee rijen: $190 + 392 = 582$. Een positieve test ontstaat langs twee wegen: besmet en positief, of niet besmet en positief. Zo vind je ook de kans op een gebeurtenis $B$ die langs twee paden kan ontstaan.

\begin{theorie}
\begin{stelling}{Wet van de totale kans}
Voor gebeurtenissen $A$ en $B$ met $0 < P(A) < 1$ is

\vspace{1mm}
\qquad $\formule{P(B) = P(A) \cdot P(B \gegeven A) + P(\compl{A}) \cdot P(B \gegeven \compl{A})}$
\end{stelling}

\begin{bewijs}
$B$ valt uiteen in de delen $A \cap B$ en $\compl{A} \cap B$, en die zijn onverenigbaar. De somregel en daarna de productregel op elk deel geven
\bewijsregel{P(B) = P(A \cap B) + P(\compl{A} \cap B) = P(A) \cdot P(B \gegeven A) + P(\compl{A}) \cdot P(B \gegeven \compl{A}) .}
\end{bewijs}
\end{theorie}

In een boomdiagram is dat de som van de twee paden die in $B$ eindigen. Valt $\Omega$ in meer dan twee delen uiteen, zoals drie machines in een fabriek, dan zijn er meer paden en dus meer termen. Het bewijs blijft hetzelfde.

\voorbeeld{}{Bereken met de gegeven percentages, zonder tabel, de kans dat een willekeurige inwoner positief test.\label{vb:h2-totaal}}{%
\uitleg{Een positieve test komt langs twee paden: besmet ($A$) of niet besmet ($\compl{A}$). We vermenigvuldigen langs elk pad en tellen op.}
\naast[4.6cm]{\begin{stappenvrij}{5.6cm}{4cm}
P(A) = 0{,}02, \ P(\compl{A}) = 0{,}98 & De eerste vertakking. \\[1ex]
P(+ \gegeven A) = 0{,}95, \ P(+ \gegeven \compl{A}) = 0{,}04 & De tweede vertakking. \\[1ex]
P(+) = 0{,}02 \cdot 0{,}95 + 0{,}98 \cdot 0{,}04 & Wet van de totale kans. \\[1ex]
= 0{,}019 + 0{,}0392 = 0{,}0582 & Dat is $\frac{582}{10\,000}$, zoals in de tabel.
\end{stappenvrij}}{\begin{tikzpicture}[x=1.9cm,y=.55cm]
  \node[blad] (s) at (0,0) {};
  \node[blad] (a) at (1,1.6) {$A$}; \node[blad] (na) at (1,-1.6) {$\compl{A}$};
  \draw[tak] (s) -- node[kanslabel,above left] {$0{,}02$} (a); \draw[tak] (s) -- node[kanslabel,below left] {$0{,}98$} (na);
  \node[blad] (ap) at (2,2.4) {$+$}; \node[blad] (am) at (2,.8) {$-$};
  \node[blad] (nap) at (2,-.8) {$+$}; \node[blad] (nam) at (2,-2.4) {$-$};
  \draw[tak] (a) -- node[kanslabel,above] {$0{,}95$} (ap); \draw[tak] (a) -- node[kanslabel,below] {$0{,}05$} (am);
  \draw[tak] (na) -- node[kanslabel,above] {$0{,}04$} (nap); \draw[tak] (na) -- node[kanslabel,below] {$0{,}96$} (nam);
\end{tikzpicture}}}

\begin{oefeningen}
\oef{h2-totaal}
\end{oefeningen}

\needspace{16\baselineskip}
\subsection{De wet van Bayes}

Een labo meet hoe vaak de test positief is bij besmette mensen: $P(\text{positief} \gegeven \text{besmet})$. De patiënt wil het omgekeerde weten: $P(\text{besmet} \gegeven \text{positief})$. De wet van Bayes keert een voorwaardelijke kans om.

\begin{theorie}
\begin{stelling}{Wet van Bayes}
Voor gebeurtenissen $A$ en $B$ met $P(A) > 0$ en $P(B) > 0$ is

\vspace{1mm}
\qquad $\formule{P(A \gegeven B) = \dfrac{P(A) \cdot P(B \gegeven A)}{P(B)}}$
\end{stelling}

\begin{bewijs}
Volgens de definitie is $P(A \gegeven B) = \frac{P(A \cap B)}{P(B)}$, en volgens de productregel is $P(A \cap B) = P(A) \cdot P(B \gegeven A)$. Vul dat in in de teller.
\end{bewijs}
\end{theorie}

In een boomdiagram is de teller het ene pad dat je zoekt, en de noemer de som van alle paden die in $B$ eindigen: de wet van de totale kans. Bij de sneltest is dat $\frac{0{,}019}{0{,}0582} \approx 0{,}33$, hetzelfde als met de tabel.

\voorbeeld{}{Van alle mails is 30~\% spam. Het woord \textit{gratis} staat in 40~\% van de spammails en in 2~\% van de gewone mails. Een mail bevat \textit{gratis}. Wat is de kans dat het spam is?\label{vb:h2-bayes-formule}}{%
\uitleg{Gegeven zijn kansen op \textit{gratis} als je weet of een mail spam is. Gevraagd is het omgekeerde: Bayes, met de totale kans in de noemer.}
\begin{stappenbreed}
P(G) = 0{,}3 \cdot 0{,}4 + 0{,}7 \cdot 0{,}02 = 0{,}12 + 0{,}014 = 0{,}134 & Totale kans op \textit{gratis}: via spam of via een gewone mail. \\[1ex]
P(S \gegeven G) = \frac{0{,}3 \cdot 0{,}4}{0{,}134} = \frac{0{,}12}{0{,}134} \approx 0{,}90 & Bayes: het pad spam en \textit{gratis}, gedeeld door alle paden naar \textit{gratis}.
\end{stappenbreed}
Het woord \textit{gratis} verhoogt de kans op spam van 30~\% naar ongeveer 90~\%.}

\begin{werkwijze}
Maak een tabel met aantallen voor een ronde groep (bv. 10\,000), of een boomdiagram met de gegeven kansen. Zoek de totale kans op wat je waarneemt (zoals \textit{positief}) door alle paden ernaar op te tellen. De gevraagde voorwaardelijke kans is het gunstige pad gedeeld door die totale kans.
\end{werkwijze}

\begin{oefeningen}
\oef{h2-bayes-tabel}

\oef{h2-bayes}
\end{oefeningen}

%========================================================================================
\opdrachtenkop

\begin{oefeningen}
\needspace{18\baselineskip}
\oef{h2-o13}

\oef{h2-o2}

\oef{h2-o16}

\oef{h2-o1}

\needspace{10\baselineskip}
\oef{h2-o14}

\oef{h2-o4}

\oef{h2-o17}

\oef{h2-o3}

\oef{h2-o15}

\oef{h2-o5}

\needspace{12\baselineskip}
\oef{h2-o23}

\oef{h2-o24}

\oef{h2-o19}

\oef{h2-o7}

\oef{h2-o21}

\needspace{14\baselineskip}
\oef{h2-o20}

\needspace{10\baselineskip}
\oef{h2-o6}

\oef{h2-o22}

\needspace{10\baselineskip}
\oef{h2-o18}

\oef{h2-o8}

\oef{h2-o9}

\oef{h2-o11}

\oef{h2-o25}

\oef{h2-o26}

\oef{h2-o10}

\oef{h2-o12}
\end{oefeningen}

%----------------------------------------------------------------------------------------
% KOPPEL: verbanden met andere leerstof
%----------------------------------------------------------------------------------------
\needspace{20\baselineskip}
\opdrachtsectie{Koppel}

\lkr{De oefening gebruikt een overgangsmatrix uit Matrices (hoofdstuk Matrixmodellen).}

\begin{oefeningen}

\oef{h2-koppel-2}
\end{oefeningen}

%----------------------------------------------------------------------------------------
\opdrachtsectie{Test jezelf}

\begin{oefeningen}
\oef{h2-t1}

\oef{h2-t2}

\oef{h2-t3}

\oef{h2-t4}

\oef{h2-t5}

\oef{h2-t6}

\oef{h2-t7}
\end{oefeningen}

\herhaalschema{%
  \fv{oef:h2-t1}{vb:h2-kruistabel}{oef:h2-kruistabel}%
  \fv{oef:h2-t2}{vb:h2-boom}{oef:h2-boom}%
  \fv{oef:h2-t3}{vb:h2-onafhankelijk}{oef:h2-onafhankelijk}%
  \fv{oef:h2-t4}{vb:h2-bayes}{oef:h2-bayes-tabel}%
  \fv{oef:h2-t5}{vb:h2-herhaling}{oef:h2-onafhankelijk-2}%
  \fv{oef:h2-t6}{vb:h2-bayes-formule}{oef:h2-bayes}%
  \fv{oef:h2-t7}{vb:h2-onverenigbaar}{oef:h2-onverenigbaar}}

\ifzelfstudie
\section*{Maak je eigen samenvatting}
Sluit dit boek en schrijf in je schrift: wat een voorwaardelijke kans is en hoe je ze in een kruistabel afleest, het verschil tussen $P(A \gegeven B)$ en $P(B \gegeven A)$, de productregel en hoe je een boomdiagram gebruikt, wanneer gebeurtenissen onafhankelijk zijn (en het verschil met onverenigbaar), de wet van de totale kans en de wet van Bayes. Vergelijk daarna met de kaders.
\fi
